% Parked from 02-v2 main.tex on 2026-10-31 (author's decision: core-only version, option B).
% Contents: Theorem C(i) -- polynomial data for the binary typing onto U_1 at fan-out two --
% with its full proof (old Appendix G); the comparison appendix with Yoshinaka's
% multidimensional substitutability and the L_R example; the self-contained substitutability
% proof of the bracketed doubling family (now cited from the companion, Prop. 6.4(i), App. C);
% the full proofs of Theorem poly and Corollary gold-explanatory; the earlier versions of the
% introduction and conclusion paragraphs that mentioned these results.
% All statements were proved and externally checked while in the manuscript (README change log,
% reviews/). Together with parked_u1_typing_level.tex this is the seed of a separate paper on
% typings onto U_1 across fan-out (Open Problems S1 and fanout-two).
% This file is not compiled; labels are kept unchanged so that blocks can be restored.

%==============================================================================
% Section 8.2: Theorem binary-u1-polydata (binary typing onto U1 at fan-out two), outline and scope remark
%==============================================================================
\subsection{Polynomial data for the single-neutral-letter typing}
\label{subsec:binary-u1}

Let $\Sigma=\{\mathtt a,\mathtt b\}$ and let $s_{U,2}:\Sigma^+\twoheadrightarrow U_1$
send $\mathtt a$ to $1$ and $\mathtt b$ to $\zeta$; a word has type $\zeta$ exactly
when it contains $\mathtt b$.  Unlike the typing $s_U$ of
Section~\ref{sec:fanout-one}, it has a single letter of type $1$.  At fan-out two
the arity-two condition gives a polynomial bound.

\begin{theorem}[Polynomial data for the binary typing onto $U_1$]
\label{thm:binary-u1-polydata}
Let $r_{\max}\ge2$ and put
$C_{r_{\max}}:=42r_{\max}(r_{\max}+1)+96(r_{\max}+1)$.  For every
$L\in\Cmcf{2}{s_{U,2}}$ and every reduced good fan-out-$\le2$,
rank-$\le r_{\max}$ presentation $G$ of $L\cap\Sigma^+$, the operators
$\mathcal B^{\le r_{\max}}_{2,s_{U,2}}$ and $\mathcal B_{2,s_{U,2}}$ have a common
characteristic sample $K$ with
\[
 \|K\|_+\le 3^{2r_{\max}}\bigl(3+r_{\max}C_{r_{\max}}\bigr)\,
 \|G\|^3\bar\theta_G^{\,5}+\mathbf 1_{\{\eps\in L\}}.
\]
In particular, the first question of Open Problem~\ref{op:fanout-two} has a
positive answer at $f=2$ for every $r_{\max}\ge2$, although
$U_1\notin\mathbf{LG}$.
\end{theorem}

\begin{proof}[Outline]
Say that $L$ has a \emph{deep insertion} if some word of $L$ has a factor
$\mathtt a^6z\mathtt a^6$ with $z$ containing $\mathtt b$.  Without one, marking the
first and last six letters of every pure root component keeps the type of
every shortened tuple and bounds the representative length $\tau$ by
$O(\|G\|\bar\theta_G)$ (Lemma~\ref{lem:app-u1-shallow}).  With one, two-hole
contexts around the run $\mathtt a^6$ identify $(\mathtt a^i,\mathtt a^j)$ and
$(\mathtt a^{i'},\mathtt a^{j'})$ with $i+j=i'+j'$, which forces a factor
$\mathtt a\mathtt b\mathtt a\mathtt a$ and then $\D(\mathtt{ab})=\D(\mathtt{ba})$, so
$L$ is Parikh-saturated (Lemma~\ref{lem:app-u1-deep}).  For Parikh-saturated
targets an analysis of node-minimal derivations bounds $\tau$ by
$O(\|G\|^2\bar\theta_G^5)$ (Lemma~\ref{lem:app-u1-parikh}), and
Corollary~\ref{cor:distribution-representatives} converts the bound on $\tau$
into the sample bound.  The details are in Appendix~\ref{app:binary-u1}.
\end{proof}

\begin{remark}[Scope of Theorem~\ref{thm:binary-u1-polydata}]
\label{rem:binary-u1-scope}
(a) The theorem concerns one typing and fan-out two.  It decides neither the
pseudovariety question at fan-out two (Open Problem~\ref{op:fanout-two}) nor
whether $s_{U,2}$ needs superpolynomial data at fan-out one (Open
Problem~\ref{op:S1}).
(b) The proof uses the arity-two condition only through runs of the single
neutral letter; with two neutral letters a long pure factor need not contain a
long run.  For every typing onto $U_1$ and every fan-out,
Lemma~\ref{lem:app-u1-shallow} bounds $\tau(G,s)$ linearly in the flank length
of deep insertions, and Lemma~\ref{lem:app-u1-runs} bounds
$t_{\mathrm{typ}}(G,s)\ge\tau(G,s)$ linearly in the number of maximal
factors without letters of type $1$, without using substitutability.  So any
superpolynomial family for such a typing must contain deep insertions whose
pure flanks grow superpolynomially and words with superpolynomially many such
factors.
(c) The proof treats tuples of arity at most two; fan-out $f\ge3$ is open.
\end{remark}

%==============================================================================
% Appendix G: polynomial data for the binary typing onto U1 at fan-out two (proof of Theorem binary-u1-polydata)
%==============================================================================
\section{Polynomial Data for the Binary Typing onto \texorpdfstring{$U_1$}{U1}}
\label{app:binary-u1}

This appendix proves Theorem~\ref{thm:binary-u1-polydata}.  Throughout,
$G=(V,\Sigma,P,S,\fan)$ is reduced and good of fan-out $\le f$ and rank
$\le r$, $r:=r_{\max}$, $v:=|V|$, $\theta:=\bar\theta_G$, and for each
$A\in V$ the context $E_A$ is that of Lemma~\ref{lem:thickness-equivalence}(iv),
so $|E_A|_{\mathrm{fix}}\le\theta$ and $E_A[L_A(G)]\subseteq L(G)$.  For each $B\in V$
fix $\omega_B\in L_B(G)$ with $\|\omega_B\|_1\le\theta$.  In a derivation tree, the
\emph{yield} of a node is the tuple it derives, and a node \emph{writes} the
terminals that occur in its template.

\begin{lemma}\label{lem:app-u1-fillers}
Let $L\in\Cmcf{f}{s}$ with $L(G)=L\cap\Sigma^+$, and let $X=A_{\mathbf t}$ be a typed
nonterminal of arity $d$.  Then
$\operatorname{Cl}_{L,s}(X)=\{\tuple z\in(\Sigma^+)^d:s^{(d)}(\tuple z)=\mathbf t,\
E_A[\tuple z]\in L\}$; hence $\tau_X$ is the least total length of such a tuple.
\end{lemma}

\begin{proof}
Fix $\tuple u\in L_X$.  By Lemma~\ref{lem:typed-homogeneity},
$\operatorname{Cl}_{L,s}(X)$ is the set of tuples with the type and the
distribution of $\tuple u$; they are accepted by $E_A\in\D_L(\tuple u)$.  A tuple
$\tuple z$ of type $\mathbf t$ accepted by $E_A$ shares this context with $\tuple u$,
so $\D_L(\tuple z)=\D_L(\tuple u)$ by substitutability.
\end{proof}

\begin{lemma}\label{lem:app-u1-identity}
In a derivation tree of $G$, let $u$ be a proper ancestor of $u'$ with the same
nonterminal, and suppose their yields have the same total Parikh vector.  Then
every rule on the path from $u$ to $u'$ has rank one, writes no terminal and
is a componentwise identity, so replacing the subtree at $u$ by the subtree at
$u'$ changes no derived tuple and removes at least one node.
\end{lemma}

\begin{proof}
By linearity and nondeletion every letter of the yield of $u'$ occurs in the
yield of $u$, so equal Parikh vectors exclude terminals written on the path
and siblings off the path, since siblings have nonempty yields
($\lambda$-freeness).  So every rule on the path has one child of some fan-out
$d'$, each of its $d'$ variables occurs once (nondeletion) and in order
(nonpermutation), and each parent component is nonempty ($\lambda$-freeness).
Two variables in one component would be consecutive, which nonmerging forbids;
hence every component is a single variable, the fan-outs agree, and the rule is
the identity.
\end{proof}

For a typing $s$ onto $U_1$ put $\Sigma_1:=\{x\in\Sigma:s(x)=1\}$; words of
$\Sigma_1^+$ are \emph{pure}, and a word has type $\zeta$ iff it is not pure.  For
$c\ge1$, $L$ has a \emph{$c$-deep insertion} if some word of $L$ has a factor
$xzy$ with $x,y\in\Sigma_1^c$ and $z$ not pure.

\begin{lemma}\label{lem:app-u1-shallow}
Let $s$ be a typing onto $U_1$, $c\ge1$, and suppose $L(G)=L\cap\Sigma^+$ has no
$c$-deep insertion.  Then $\tau(G,s)\le2^{f+1}cf(r+1)v\theta$.
\end{lemma}

\begin{proof}
Let $X=A_{\mathbf t}$ be a typed nonterminal of arity $d$ and $T$ a derivation of
$\widetilde G_0$ from $X$ with the fewest nodes.  No typed label repeats on a path
of $T$, since cutting between two equal typed labels gives a smaller derivation
from $X$; as all components are nonempty, there are at most $2^fv$ typed labels,
which bounds the length of every path.  In each pure root component mark the
first and the last $\min\{c,|u_i|\}$ letters, and in each other component one
letter outside $\Sigma_1$; there are at most $2cf$ marks.  Let $H$ be the set of
nodes on the paths from the root to the nodes writing marked letters, so
$|H|\le2cf\cdot2^fv$.  Erase the type indices and replace, one at a time, each
subtree hanging off $H$ by a derivation of $\omega_B$ for its root nonterminal $B$.

Every intermediate root tuple lies in $L_A(G)$.  We claim that the pure root
components stay pure.  A replaced subtree contributes contiguous blocks to root
components, and these blocks contain no mark.  A pure component of length at
most $2c$ is entirely marked and receives no block; in a longer one every block
lies between the first and the last $c$ letters, which are unchanged and pure.
If a replacement made such a component impure, inserting the root tuple into
$E_A$ would give a word of $L$ with a $c$-deep insertion.  Marked letters outside
$\Sigma_1$ also stay, so the final tuple has type $\mathbf t$ and lies in $L_X$ by
Lemma~\ref{lem:witness-length-template}(i).  Each node of $H$ writes at most
$\theta$ letters and has at most $r$ replaced children of length at most
$\theta$, so the tuple has length at most $(r+1)\theta|H|$.
\end{proof}

The next lemma needs no substitutability.  A \emph{$\zeta$-run} of a word is a
maximal factor that contains no letter of $\Sigma_1$.

\begin{lemma}\label{lem:app-u1-runs}
Let $s$ be a typing onto $U_1$ and $k\ge0$, and suppose every word of $L(G)$ has
at most $k$ $\zeta$-runs.  Then
$t_{\mathrm{typ}}(G,s)\le2^f(r+1)\bigl(2f+(2f+1)k\bigr)v\theta$.
\end{lemma}

\begin{proof}
Let $X=A_{\mathbf t}$ be a typed nonterminal and $T$ a derivation of
$\widetilde G_0$ from $X$ whose yield $\tuple u$ has the least total length and,
among those, the fewest nodes.  Cutting between two equal typed labels on a path
gives a derivation from $X$ with a yield that is not longer and with fewer
nodes, so paths of $T$ have at most $2^fv$ nodes.  Mark the first and the last
letter of each pure root component and one letter outside $\Sigma_1$ in each
other component; there are at most $2f$ marks.  Call a node heavy if its yield
has total length greater than $\theta$.

Let $Q$ be an antichain of heavy nodes and $Q_0\subseteq Q$ the nodes whose
subtrees write no mark, so $|Q\setminus Q_0|\le2f$.  For $h\in Q_0$ with
nonterminal $C_h$, erase the type indices and replace the subtree at $h$ by a
derivation of $\omega_{C_h}$.  The new root tuple lies in $L_A(G)$ and is shorter,
since by nondeletion every letter of the yield of $h$ occurs in $\tuple u$.  The
marks keep every impure root component impure.  If every pure root component
stayed pure, the new tuple would have type $\mathbf t$ and lie in $L_X$ by
Lemma~\ref{lem:witness-length-template}(i), contradicting the choice of $T$.
Hence $\omega_{C_h}$ puts a letter outside $\Sigma_1$ into a block of $h$ lying in
an originally pure root component.

Each $h\in Q_0$ owns at most $f$ blocks, which are disjoint intervals of root
components and contain no mark.  Join two owners when blocks of theirs in a pure
root component are adjacent with no letter between them.  Each block has at most
two adjacent blocks, so this graph has maximum degree at most $2f$ and an
independent set $I$ with $|I|\ge|Q_0|/(2f+1)$.  Replace the subtrees at all
$h\in I$ at once, giving $\tuple u'\in L_A(G)$.  In a pure root component, two
blocks of distinct owners in $I$ are separated by an unchanged letter of
$\Sigma_1$: either their gap was nonempty, or it contains a block of an owner
outside $I$, which is kept.  The first and the last letter of the component are
marks and stay.  Hence every $h\in I$ contributes a $\zeta$-run of
$E_A[\tuple u']\in L(G)$, different owners contribute different runs, so
$|I|\le k$, $|Q_0|\le(2f+1)k$ and $|Q|\le2f+(2f+1)k$.

The heavy nodes are closed under ancestors.  Those without heavy children form
an antichain, and every heavy node lies on the path from the root to one of
them; so there are at most $2^fv(2f+(2f+1)k)$ heavy nodes.  Every letter of
$\tuple u$ is written by a heavy node or lies in a light child of one, and each
heavy node writes at most $\theta$ letters and has at most $r$ children.  Hence
$\|\tuple u\|_1\le\max\{\theta,(r+1)\theta\cdot2^fv(2f+(2f+1)k)\}$.
\end{proof}

From now on $\Sigma=\{\mathtt a,\mathtt b\}$, $s=s_{U,2}$, $\Psi(w):=(|w|_{\mathtt a},|w|_{\mathtt b})$,
and $L$ is \emph{Parikh-saturated} if membership of a word of $\Sigma^+$ in $L$
depends only on its Parikh vector.

\begin{lemma}\label{lem:app-u1-transfer}
Let $L\in\Cmcf{2}{s}$ and let some word of $L$ contain $\mathtt a^m$.
\textup{(i)} If $i,j,i',j'\ge1$ and $i+j=i'+j'\le m-1$, then
$U\mathtt a^iZ\mathtt a^jV\in L\iff U\mathtt a^{i'}Z\mathtt a^{j'}V\in L$ for all
$U,V\in\Sigma^*$ and $Z\in\Sigma^+$.
\textup{(ii)} If $m\ge4$ and some word of $L$ contains $\mathtt a\mathtt a\mathtt b\mathtt a$
or $\mathtt a\mathtt b\mathtt a\mathtt a$, then $L$ is Parikh-saturated.
\end{lemma}

\begin{proof}
(i) Let $p\mathtt a^mq\in L$ and $k:=m-i-j\ge1$.  The tuples $(\mathtt a^i,\mathtt a^j)$ and
$(\mathtt a^{i'},\mathtt a^{j'})$ of type $(1,1)$ share the admissible context
$p\blank_1\mathtt a^k\blank_2q$, so they have the same distribution; test it on the
admissible context $U\blank_1Z\blank_2V$.
(ii) Say $U\mathtt a\mathtt a\mathtt b\mathtt a V\in L$ (the other case is symmetric).  By (i)
with $i+j=3$, $U\mathtt a\mathtt b\mathtt a\mathtt aV\in L$.  The context $(U\mathtt a,\mathtt aV)$
accepts $\mathtt{ab}$ and $\mathtt{ba}$, both of type $\zeta$, so
$\D(\mathtt{ab})=\D(\mathtt{ba})$: in every context $\mathtt{ab}$ may be replaced by
$\mathtt{ba}$ and conversely.  Adjacent transpositions connect all words with one
Parikh vector.
\end{proof}

\begin{lemma}\label{lem:app-u1-deep}
Let $L\in\Cmcf{2}{s}$.  If some word of $L$ has a factor $\mathtt a^6z\mathtt a^6$ with $z$
containing $\mathtt b$, then $L$ is Parikh-saturated.
\end{lemma}

\begin{proof}
The first maximal run $\mathtt b^e$ of $z$ is preceded by at least six letters
$\mathtt a$ and followed by at least one, so the word is
$P\mathtt a^i\mathtt b^e\mathtt a^jR$ with $i\ge6$, $j\ge1$.  Choose $i_0+j_0=5$ with
$1\le i_0\le i$, $1\le j_0\le j$, and put $P_0:=P\mathtt a^{i-i_0}$,
$R_0:=\mathtt a^{j-j_0}R$.  Lemma~\ref{lem:app-u1-transfer}(i) with $m=6$ gives
$w_k:=P_0\mathtt a^k\mathtt b^e\mathtt a^{5-k}R_0\in L$ for $1\le k\le4$.  If $e=1$, $w_2$
contains $\mathtt{aaba}$.  If $e\ge2$, the tuples $\tuple x=(\mathtt a,\mathtt b^{e-1}\mathtt a)$
and $\tuple y=(\mathtt a\mathtt a,\mathtt b^{e-1})$ of type $(1,\zeta)$ are accepted by
$P_0\blank_1\mathtt b\blank_2\mathtt a^3R_0$ (giving $w_1$ and $w_2$), so they have the
same distribution; the context $P_0\blank_1\mathtt a\blank_2\mathtt b\mathtt a^2R_0$ accepts
$\tuple y$ (giving $w_3$) and hence $\tuple x$, which gives a word containing
$\mathtt{abaa}$.  Apply Lemma~\ref{lem:app-u1-transfer}(ii).
\end{proof}

\begin{lemma}\label{lem:app-u1-parikh}
Let $f\le2$, let $L\in\Cmcf{f}{s}$ be Parikh-saturated, and
$L(G)=L\cap\Sigma^+$.  Then $\tau(G,s)\le(42r(r+1)+12(r+1))v^2\theta^5$.
\end{lemma}

\begin{proof}
Fix $X=A_{\mathbf t}$, put $(p,q):=\Psi(E_A)$, so $p+q\le\theta$, and
$\mathcal P:=\Psi(L\cap\Sigma^+)$.  By Lemma~\ref{lem:app-u1-fillers} and Parikh
saturation, $\tau_X$ is the least $\|\tuple z\|_1$ over tuples $\tuple z$ of type
$\mathbf t$ with $(p,q)+\sum_i\Psi(z_i)\in\mathcal P$.  Write $(x_B,y_B):=\Psi(\omega_B)$.
For a node $u$ of a derivation tree let $(a_u,b_u)$ be the Parikh vector of its
yield and $B_u$ its nonterminal.  If $u$ lies in a successful derivation of a word
$w$, replacing its subtree by a derivation of $\omega_{B_u}$ gives a word of $L$
with Parikh vector $\Psi(w)-(a_u,b_u)+(x_{B_u},y_{B_u})$; for a proper descendant
$u'$ with $B_{u'}=B_u$, replacing the subtree at $u$ by that at $u'$ gives one
with $\Psi(w)-(a_u-a_{u'},b_u-b_{u'})$.

\emph{Type $\zeta$ only.}  Mark one letter $\mathtt b$ in each component and argue as
in Lemma~\ref{lem:app-u1-shallow}: a derivation from $X$ with the fewest nodes repeats
no typed label on a path, so its paths have at most $2^fv\le4v$ nodes; the paths to
the at most two marks contain at most $8v$ nodes, and replacing every subtree hanging
off them by a derivation of the corresponding $\omega_B$ keeps the marks.  No
component is pure, so no hypothesis on deep insertions is needed, and the marks keep
every component of type $\zeta$; hence $\tau_X\le2\cdot4v(r+1)\theta$.

\emph{Mixed type}, say $\mathbf t=(1,\zeta)$ (the other order is symmetric).  The
fillers may be taken as $(\mathtt a^c,\mathtt b^y)$, so $\tau_X$ is the least $c+y$
with $c,y\ge1$ and $(p+c,q+y)\in\mathcal P$.  Fix such $c,y$ and a successful
derivation $T$ of $W:=\mathtt b^{q+y}\mathtt a^{p+c}$ with the fewest nodes.  Call $u$
heavy if $a_u+b_u>\theta$.  If $a_u<c$ and $b_u<y$ for a heavy $u$, its replacement
gives $(p+c',q+y')\in\mathcal P$ with $c',y'\ge1$ and $c'+y'<c+y$; so every heavy
node has $a_u\ge c$ or $b_u\ge y$.  If $u'$ is a proper descendant of $u$ with the
same nonterminal, then $a_u-a_{u'}\ge c$ or $b_u-b_{u'}\ge y$: a zero difference
contradicts node-minimality (Lemma~\ref{lem:app-u1-identity}), and another small
difference contradicts minimality of $c+y$.  Along a path, $(\lfloor a_u/c\rfloor,
\lfloor b_u/y\rfloor)$ is componentwise nonincreasing and takes at most $p+q+3\le
\theta+3$ values, and no nonterminal repeats while it is constant; so paths have
at most $v(\theta+3)$ nodes.  Heavy nodes at one depth have disjoint yields, so
there are at most $(p+c)/c+(q+y)/y\le\theta+2$ of them.  Heavy nodes are closed
under ancestors, and every letter is written by a heavy node or lies in one of its
at most $r$ non-heavy children; hence
$c+y\le|W|\le\max\{\theta,(r+1)\theta(\theta+2)(\theta+3)v\}\le12(r+1)v\theta^3$.

\emph{Type $1$ only}, with $h:=d\in\{1,2\}$.  Let $K:=\min\{k\ge h:(p+k,q)\in\mathcal P\}$;
then $\tau_X=K$, and we may assume $K>\theta+2\ge p+2$.  Fix a successful derivation
$T$ of $W:=\mathtt b^q\mathtt a^{p+K}$ with the fewest nodes.

(P1) \emph{Row rigidity}: there are no $(c_0,y),(c_0+\delta,y)\in\mathcal P$ with
$c_0\ge1$ and $1\le\delta\le K-h$.  Otherwise the context
$(\mathtt b^y\mathtt a^{c_0-1},\eps)$ accepts $\mathtt a$ and $\mathtt a^{1+\delta}$, so
$\D(\mathtt a)=\D(\mathtt a^{1+\delta})$ by $(1,s)$-substitutability; the context
$(\mathtt b^q\mathtt a^{p+K-\delta-1},\eps)$ accepts $\mathtt a^{1+\delta}$, hence $\mathtt a$, and
$(p+K-\delta,q)\in\mathcal P$ with $h\le K-\delta<K$, contradicting the choice of $K$.

(P2) \emph{Height.}  If $u'$ is a proper descendant of $u$ with the same nonterminal
and $b_u=b_{u'}$, then $a_u-a_{u'}\ge K-h+1$: the value $0$ contradicts
node-minimality (Lemma~\ref{lem:app-u1-identity}), and a value in $[1,K-h]$
contradicts the choice of $K$.  Along a path $b_u$ takes at most $q+1$ values, and
three nodes with one nonterminal and one value of $b_u$ would give
$p+K\ge a_u\ge2(K-h+1)$, false as $K>p+2$.  So paths have at most $2v(q+1)$ nodes.
A node is \emph{mixed} if its yield contains $\mathtt b$; each mixed node lies on the
path to one of the $q$ nodes writing a letter $\mathtt b$, so there are
$M\le2vq(q+1)\le4v\theta^2$ of them.

(P3) \emph{Heavy pure nodes.}  Call a pure node $u$ heavy if $\ell_u:=a_u>\theta+1$,
and put $g_u:=\ell_u-x_{B_u}\ge2$.  Two comparable heavy pure nodes have different
nonterminals, since by (P2) the upper one would have
$\ell_u\ge K-h+1+\theta+2>p+K=|W|_{\mathtt a}$.  In particular heavy pure nodes with one
nonterminal have disjoint yields.  For an antichain $J$ of heavy pure nodes,
replacing all of them gives a word with Parikh vector
$(p+K-\sigma(J),q+Y(J))$, where $\sigma(J):=\sum_{u\in J}g_u$ and
$Y(J):=\sum_{u\in J}y_{B_u}$; by (P1), two antichains with
$\sigma(J),\sigma(J')\le K-h$ and $Y(J)=Y(J')$ have $\sigma(J)=\sigma(J')$, because
both $\mathtt a$-counts are at least $p+h\ge1$.

(P4) \emph{Repeated nonterminals.}  If $B$ labels $m_B\ge2$ heavy pure nodes, their
yields are disjoint, so each has $\ell_u\le p+K-(\theta+2)\le K-2$ and
$g_u\le K-2\le K-h$; comparing two singletons by (P3) gives a common value $g_B$, and
then a common $\ell_B$.  If moreover $y_B=0$, one replacement gives
$(p+K-g_B,q)\in\mathcal P$ with $h\le K-g_B<K$, which is impossible; so $y_B\ge1$.
Call $B$ \emph{frequent} if $m_B>\theta$.  For frequent $B,C$, choose $y_C$
nodes labelled $B$ and, separately, $y_B$ nodes labelled $C$, which is possible
since $m_B>\theta\ge y_C$ and $m_C>\theta\ge y_B$; both sets are antichains with
$Y=y_By_C$, and
$y_Cg_B\le(m_B-1)\ell_B\le p+K-\ell_B\le K-2\le K-h$, and likewise for $C$.  By (P3),
$y_Cg_B=y_Bg_C$, so all frequent $B$ share a ratio $\alpha:=g_B/y_B$, and
$\ell_B=x_B+\alpha y_B\le(1+\alpha)\theta$.

(P5) \emph{The ratio.}  Put $\alpha_*:=(r+1)(\theta+1)(2\theta+1)v\le6(r+1)v\theta^2$;
we show $\alpha\le\alpha_*$ when frequent nonterminals exist.  This is clear if
$\alpha\le\theta$.  Otherwise fix a heavy pure node $u$ labelled by a
frequent nonterminal $B$, and let $H_u$ be the number of heavy pure nodes in
its subtree.  These nodes are closed under ancestors within that subtree.
Every letter in the yield of $u$ is written by a heavy pure node or belongs
to the yield of a non-heavy child of such a node.  Each heavy pure node writes
at most $\theta$ letters and has at most $r$ non-heavy children, each of total
yield length at most $\theta+1$; hence
$\ell_B\le(r+1)(\theta+1)H_u$.  Group them by their number of heavy proper
ancestors in the subtree: by (P3) there are at most $v$ groups, each an
antichain.  In one group the nodes with frequent labels have length at least
$\alpha$ and disjoint yields inside that of $u$, so there are at most
$(1+\alpha)\theta/\alpha<\theta+1$ of them; nodes with other labels number at
most $v\theta$ in all of $T$.  So $H_u\le v(2\theta+1)$, and
$\alpha\le g_B\le\ell_B\le\alpha_*$.

(P6) \emph{Counting.}  Call a pure node \emph{large} if $\ell_u>T_*:=(\alpha_*+1)\theta
\le7(r+1)v\theta^3$.  By (P4)--(P5) no frequent nonterminal labels a large node,
so there are at most $v\theta$ large nodes, and they are closed under ancestors
within pure subtrees.  Every letter $\mathtt a$ of $W$ is written by a mixed node, is
written by a large node, or lies in a maximal non-large pure subtree, whose root
is the root of $T$ or a child of a mixed or a large node.  Hence
\[
 p+K\le\theta M+(rM+rv\theta+1)T_*+\theta\cdot v\theta
 \le5v\theta^3+6rv\theta^2\cdot7(r+1)v\theta^3\le(42r(r+1)+5)v^2\theta^5,
\]
using $M\le4v\theta^2$ and $4rv\theta^2+rv\theta+1\le6rv\theta^2$.  In all three
cases $\tau_X\le(42r(r+1)+12(r+1))v^2\theta^5$.
\end{proof}

\begin{proof}[Proof of Theorem~\ref{thm:binary-u1-polydata}]
If $L$ has no $6$-deep insertion, Lemma~\ref{lem:app-u1-shallow} with $f=2$,
$c=6$ gives $\tau\le96(r+1)v\theta$.  Otherwise $L$ is Parikh-saturated by
Lemma~\ref{lem:app-u1-deep}, and Lemma~\ref{lem:app-u1-parikh} applies.  In both
cases $\tau\le C_rv^2\theta^5$.  By Corollary~\ref{cor:distribution-representatives}
there is a common characteristic sample with
$\|K\|_+\le N_{\mathrm{typ}}(1+2\theta+r\tau)+\mathbf 1_{\{\eps\in L\}}$, and
$N_{\mathrm{typ}}\le|V||\mathsf S^\sharp|^2+|P||\mathsf S^\sharp|^{2r}\le3^{2r}\|G\|$ since
$|\mathsf S^\sharp|=3$.  With $v\le\|G\|$ and $1+2\theta\le3\|G\|^2\theta^5$ this gives
the stated bound.
\end{proof}

%==============================================================================
% Appendix I: comparison with multidimensional substitutability (Yoshinaka baseline, L_R example)
%==============================================================================
\section{Comparison with Multidimensional Substitutability}
\label{sec:comparison}

Yoshinaka's multidimensional substitutability~\cite{Yoshinaka2009,Yoshinaka2011}
is the direct positive-data baseline.  This appendix shows that trivial typing
recovers it and window typing recovers the CFG $(k,\ell)$ hierarchy
of~\cite{Yoshinaka2008}, while a fixed finite typing admits non-context-free
targets outside the baseline.
Since every $(f,s)$ condition includes arity one,
$\Cmcf{f}{s}\cap\mathrm{CFL}\subseteq\Cmcf{1}{s}$.


\begin{definition}[Multidimensional substitutability]
\label{def:yoshinaka-pd-substitutability}
A language $L$ is \emph{$f$D-substitutable} if for every $1\le d\le f$, all
$E,F\in\mathcal E_d$ and all $\tuple x,\tuple y\in(\Sigma^+)^d$,
\(E[\tuple x],E[\tuple y],F[\tuple x]\in L\Rightarrow F[\tuple y]\in L\)
(the \emph{substitution implication}); $\YSub{f}$ is the class of these languages.
This is Yoshinaka's definition, with nonempty components and nonempty internal
separators \cite[Section~3.1 and footnote~3]{Yoshinaka2011}.
\end{definition}

The substitution implication is $(f,s_{\mathrm{triv}})$-tuple substitutability,
so $\Cmcf{f}{s_{\mathrm{triv}}}=\YSub{f}\cap f\text{-}\mathrm{MCFL}\subseteq\Cmcf{f}{s}$
for every $s$.  Yoshinaka \cite[Definition~2]{Yoshinaka2008} calls $L$
\emph{$k,l$-substitutable} if for all $x_1,y_1,z_1,x_2,y_2,z_2\in\Sigma^*$,
$v\in\Sigma^k$ and $u\in\Sigma^\ell$ with $vy_1u,vy_2u\ne\eps$,
$x_1vy_1uz_1,x_1vy_2uz_1,x_2vy_1uz_2\in L$ implies $x_2vy_2uz_2\in L$.  Here
$y_1,y_2$ may be empty; only the factors $vy_iu$ must be nonempty.  Exchanging
$y_1$ and $y_2$ gives the converse implication, so the condition says that two
such factors with a common context have equal distributions.  Two nonempty words
have equal $s_{k,\ell}$-type iff they are identical or both have length at least
$\max\{1,k+\ell\}$ with the same first $k$ and last $\ell$ letters, i.e.\ iff they
are of the forms $vy_1u,vy_2u$ above, with $y_1,y_2$ possibly empty, or equal.
Words shorter than $k+\ell$ are therefore compared only with themselves, in line
with Yoshinaka's remark that the condition says nothing about them.  Hence
$(1,s_{k,\ell})$-tuple substitutability is exactly Yoshinaka's $(k,\ell)$
condition; the componentwise version at higher fan-out is a definition of this
paper.

A language $L$ is \emph{strictly 2-local} if there is a set $\Pi$ of pairs over
$\Sigma\cup\{\vdash,\dashv\}$ such that a word $w$ lies in $L$ iff every pair of
adjacent symbols of $\vdash w\dashv$ lies in $\Pi$.

\begin{lemma}[Strictly 2-local languages are boundary typed]
\label{lem:sl2-boundary-control}
Every strictly 2-local $L$ is $(f,s_{1,1})$-tuple-substitutable for every $f$; in
particular $L\in\Cmcf{f}{s_{1,1}}$, and $L\in\Cmcf{f}{s}$ whenever $s_{1,1}\preceq s$.
\end{lemma}

\begin{proof}
Let $\tuple x,\tuple y\in(\Sigma^+)^d$ have equal $s_{1,1}$-types, so
corresponding components are identical letters or share first and last letters,
and let $E_0[\tuple x],E_0[\tuple y]\in L$.  Suppose $E[\tuple x]\in L$.  Since the
internal separators of $E$ are nonempty, every adjacent pair of
$\vdash E[\tuple y]\dashv$ either lies inside some $y_i$, and then occurs in
$E_0[\tuple y]\in L$, or joins a first or last letter of some $y_i$ to a symbol of
$\vdash E\dashv$, or lies in $\vdash E\dashv$, and then occurs in
$\vdash E[\tuple x]\dashv$.  So $E[\tuple y]\in L$; symmetry and $s_{1,1}\preceq s$
finish the proof (regular languages are MCFLs).
\end{proof}

\begin{proposition}[Cross-serial dependencies with finite feature licensing]
\label{prop:feature-cross-typed}
Let $\Sigma_{\mathrm{feat}}=\{\mathtt p_0,\mathtt p_1,\mathtt q_0,\mathtt q_1,
\mathtt a,\mathtt b,\mathtt c,\mathtt d\}$, let $g:\Sigma_{\mathrm{feat}}^*\to\mathbb Z_2^2$
map $\mathtt p_1\mapsto(1,0)$, $\mathtt q_1\mapsto(0,1)$ and all other letters to $0$, and
let $s:=(s_{1,1},g)|_{\Sigma_{\mathrm{feat}}^+}$ (corestricted to its image).  For
$R\subseteq\{0,1\}^2$ put
\[
 L_R:=\{\mathtt a^m\mathtt p_i\mathtt b^n\mathtt c^m\mathtt q_j\mathtt d^n:
 m,n\ge1,\ (i,j)\in R\}.
\]
Then for every $R$ and every $f\ge2$, $L_R\in\mathsf{MCFL}(2,2)\cap\Cmcf{f}{s}$, and
$L_R$ is not context-free if $R\ne\emptyset$.  For
$\mathcal R_{\mathrm{feat}}=\{(0,0),(0,1),(1,0)\}$, the language
$L_{\mathcal R_{\mathrm{feat}}}$ is neither in $\YSub{f}$ nor
$(1,s_{1,1})$-tuple-substitutable.  The typing is fixed once for the whole
family.
\end{proposition}

\begin{proof}
\emph{Grammar.}  With fan-out-two $A,B$ and rules
$A\to(\mathtt a,\mathtt c)\mid(\mathtt a\xi_1^1,\mathtt c\xi_1^2)(A)$,
$B\to(\mathtt b,\mathtt d)\mid(\mathtt b\xi_1^1,\mathtt d\xi_1^2)(B)$ and
$S\to(\xi_1^1\mathtt p_i\xi_2^1\xi_1^2\mathtt q_j\xi_2^2)(A,B)$ for
$(i,j)\in R$, we get $L_A=\{(\mathtt a^m,\mathtt c^m)\}$,
$L_B=\{(\mathtt b^n,\mathtt d^n)\}$ and $L(G)=L_R$; the grammar is good,
of fan-out and rank two.  For $(i,j)\in R$, intersecting with
$\mathtt a^+\mathtt p_i\mathtt b^+\mathtt c^+\mathtt q_j\mathtt d^+$ and erasing the
markers gives the non-context-free $\{\mathtt a^m\mathtt b^n\mathtt c^m\mathtt d^n\}$.

\emph{Typed substitutability.}  Let
$Q=\mathtt a^+\{\mathtt p_0,\mathtt p_1\}\mathtt b^+\mathtt c^+\{\mathtt q_0,\mathtt q_1\}
\mathtt d^+$ (strictly 2-local), $K_R:=\{w\in\Sigma_{\mathrm{feat}}^+:g(w)\in R\}$, and let
$\varphi$ map $\mathtt a,\mathtt c,\mathtt b,\mathtt d$ to
$(1,0),(-1,0),(0,1),(0,-1)\in\mathbb Z^2$ and the markers to $0$.  On $Q$,
$g(w)=(i,j)$ for the markers $\mathtt p_i,\mathtt q_j$ of $w$, so
$L_R=Q\cap K_R\cap\varphi^{-1}(0)$.  The three languages are
$(f,s)$-tuple-substitutable by Lemma~\ref{lem:sl2-boundary-control}, by item~(c)
after Definition~\ref{def:class-cmcf} ($K_R$ is a union of $s$-classes), and by
Proposition~\ref{prop:abelian-coset-filter}(ii);
Proposition~\ref{prop:abelian-coset-filter}(i) gives the intersection.

\emph{Failure without $g$.}  For $R=\mathcal R_{\mathrm{feat}}$, $x=\mathtt{ap_0b}$,
$y=\mathtt{ap_1b}$, $E=\blank_1\mathtt{cq_0d}$ and $F=\blank_1\mathtt{cq_1d}$ we have
$E[x],E[y],F[x]\in L_R$ but $F[y]\notin L_R$, since $(1,1)\notin R$.  Here
$s_{1,1}(x)=s_{1,1}(y)$ but $g(x)\ne g(y)$.  The same four words violate the
substitution implication of Definition~\ref{def:yoshinaka-pd-substitutability} at
$d=1$, so $L_R\notin\YSub{f}$ for every $f\ge1$.
\end{proof}

%==============================================================================
% Appendix F: self-contained substitutability proof of the bracketed doubling family (replaced by citation of the companion, Proposition 6.4(i) and Appendix C)
%==============================================================================
\section{Substitutability of the Bracketed Doubling Family}
\label{app:bracketed-family}

This appendix proves Proposition~\ref{prop:bracketed-family}(ii), adapting
the deterministic-parser argument of \cite{kuriyamaCFG}. Fix $n\ge1$, put
$L^{(0)}:=\{\mathtt a\}$ and
$L^{(i)}:=\lceil L^{(i-1)}\rceil\cup\lfloor L^{(i-1)}L^{(i-1)}\rfloor$ for
$i\ge1$, so that $L_{E_i}(G_n)=L^{(i)}$ and $L_n=L^{(n)}$. For
$x,y\in\Gamma_{\mathrm{br}}^*$ write $x\equiv y$ if
$\{(u,v):uxv\in L_n\}=\{(u,v):uyv\in L_n\}$; this is a monoid congruence.

\paragraph{The parser.}
A \emph{stack} is a word over $E_0,\ldots,E_n,\rceil,\rfloor$, with its
top on the left; $\eps$ is the empty stack. For every stack $\mathcal Z$ and
$1\le i\le n$ the moves are
\[
\begin{aligned}
E_0\mathcal Z&\xrightarrow{\mathtt a}\mathcal Z,&
E_i\mathcal Z&\xrightarrow{\lceil}E_{i-1}\rceil\mathcal Z,&
E_i\mathcal Z&\xrightarrow{\lfloor}E_{i-1}E_{i-1}\rfloor\mathcal Z,\\
\rceil\mathcal Z&\xrightarrow{\rceil}\mathcal Z,&
\rfloor\mathcal Z&\xrightarrow{\rfloor}\mathcal Z,
\end{aligned}
\]
and there are no others; in particular no move leaves $\eps$. Write
$\mathcal X\xrightarrow{w}\mathcal Y$ if reading $w$ from $\mathcal X$ ends in
$\mathcal Y$. Three properties are used.
\begin{enumerate}[label=\textup{(P\arabic*)},leftmargin=2.6em]
\item \emph{Determinism and injectivity.} For each letter $c$, the
$c$-moves form a partial injective map on stacks: an opening bracket
reveals the expanded $E_i$ and the untouched suffix, and $\mathtt a$,
$\rceil$, $\rfloor$ identify the popped top symbol. Hence every word acts
as a partial injection on stacks.
\item \emph{Weight.} Give $E_i$ weight $4^i$ and each closing marker
weight $1$, and let $\operatorname{wt}(\mathcal X)$ be the total weight. Every
move lowers $\operatorname{wt}$, since $4^i>2\cdot4^{i-1}+1$; hence
$\mathcal X\xrightarrow{w}\mathcal X$ only for $w=\eps$.
\item \emph{Locality.} Moves read and rewrite only the top symbol, and
closing markers are only popped. Hence $\mathcal X\xrightarrow{w}\mathcal Y$
implies $\mathcal X\mathcal Z\xrightarrow{w}\mathcal Y\mathcal Z$. Conversely, if
$\mathcal X_1\ne\eps$ and $\mathcal X_1\mathcal Z\xrightarrow{w}\mathcal Y$, then either
\textup{(a)} the run reaches $\mathcal Z$: $w=w_1w_2$ with
$\mathcal X_1\xrightarrow{w_1}\eps$ and $\mathcal Z\xrightarrow{w_2}\mathcal Y$, so
$\operatorname{wt}(\mathcal Y)\le\operatorname{wt}(\mathcal Z)$; or \textup{(b)} it
does not: $\mathcal Y=\mathcal Y_1\mathcal Z$ with $\mathcal X_1\xrightarrow{w}\mathcal Y_1\ne\eps$,
so $\operatorname{wt}(\mathcal Y)>\operatorname{wt}(\mathcal Z)$, and $\mathcal Y_1$
ends with the last symbol of $\mathcal X_1$ whenever that symbol is a closing
marker. For fixed $\mathcal Z$, comparing $\operatorname{wt}(\mathcal Y)$ with
$\operatorname{wt}(\mathcal Z)$ decides the case.
\end{enumerate}

\begin{lemma}\label{lem:app-br-recognition}
For $0\le i\le n$ and every stack $\mathcal Z$,
$E_i\mathcal Z\xrightarrow{w}\mathcal Z$ if and only if $w\in L^{(i)}$. Moreover, if
$\mathcal X\xrightarrow{w}\mathcal Y$ and $w\in L^{(i)}$, then $\mathcal X=E_i\mathcal Y$.
\end{lemma}

\begin{proof}
A run $E_i\mathcal Z\xrightarrow{w}\mathcal Z$ is in case~(a) of (P3), since case~(b)
ends in $\mathcal Y_1\mathcal Z\ne\mathcal Z$; so $w=w_1w_2$ with
$E_i\xrightarrow{w_1}\eps$ and $\mathcal Z\xrightarrow{w_2}\mathcal Z$, and
$w_2=\eps$ by (P2). With the lifting in (P3), it suffices to show
$E_i\xrightarrow{w}\eps$ iff $w\in L^{(i)}$, by induction on $i$. From
$E_0$ only the $\mathtt a$-move applies, and $\eps$ admits none. For
$i\ge1$, the first letter is $\lceil$ or $\lfloor$. After $\lceil$, the
stack $E_{i-1}\rceil$ must be emptied; case~(a) of (P3) with
$\mathcal X_1=E_{i-1}$, $\mathcal Z=\rceil$ splits the rest as $zz'$ with
$E_{i-1}\xrightarrow{z}\eps$ and $z'=\rceil$. After $\lfloor$, two such
splittings of $E_{i-1}E_{i-1}\rfloor$ give $z_1z_2\rfloor$ with
$E_{i-1}\xrightarrow{z_j}\eps$. Induction gives
$E_i\xrightarrow{w}\eps$ iff $w\in L^{(i)}$. For the last claim,
$E_i\mathcal Y\xrightarrow{w}\mathcal Y$ by the first part, and (P1) gives
$\mathcal X=E_i\mathcal Y$.
\end{proof}

\begin{lemma}\label{lem:app-br-exchange}
If $0\le i\le n$ and $w,w'\in L^{(i)}$, then $w\equiv w'$.
\end{lemma}

\begin{proof}
If $pwq\in L_n$, Lemma~\ref{lem:app-br-recognition} gives a run
$E_n\xrightarrow{p}\mathcal X\xrightarrow{w}\mathcal Y\xrightarrow{q}\eps$ with
$\mathcal X=E_i\mathcal Y$, and also $E_i\mathcal Y\xrightarrow{w'}\mathcal Y$; so
$pw'q\in L_n$. Exchange $w$ and $w'$ for the converse.
\end{proof}

\begin{lemma}[Branch synchronization]\label{lem:app-br-branch}
Let $\mathcal X\xrightarrow{x}\mathcal Y$ and $\mathcal X\xrightarrow{y}\mathcal Y$, where
$x,y$ are nonempty and begin with different letters. Then
$\mathcal X=E_i\mathcal Z$ for some $1\le i\le n$ and stack $\mathcal Z$, and $x=wx'$,
$y=w'y'$ with $w,w'\in L^{(i)}$, $\mathcal Z\xrightarrow{x'}\mathcal Y$ and
$\mathcal Z\xrightarrow{y'}\mathcal Y$.
\end{lemma}

\begin{proof}
Only a top $E_i$ with $i\ge1$ admits two different letters, namely
$\lceil$ and $\lfloor$. Say $x=\lceil x_0$ and $y=\lfloor y_0$, so
$E_{i-1}\rceil\mathcal Z\xrightarrow{x_0}\mathcal Y$ and
$E_{i-1}E_{i-1}\rfloor\mathcal Z\xrightarrow{y_0}\mathcal Y$. By (P3) both runs
are in the same case. Case~(b) is impossible: it would give
$\mathcal Y=\mathcal Y_1\mathcal Z=\mathcal Y_2\mathcal Z$ with $\mathcal Y_1$ ending in $\rceil$ and
$\mathcal Y_2$ in $\rfloor$. In case~(a), $x_0=x_1x'$ with
$E_{i-1}\rceil\xrightarrow{x_1}\eps$ and $\mathcal Z\xrightarrow{x'}\mathcal Y$, so
$w:=\lceil x_1$ satisfies $E_i\xrightarrow{w}\eps$ and $w\in L^{(i)}$ by
Lemma~\ref{lem:app-br-recognition}. The same argument gives $y=w'y'$.
\end{proof}

\begin{proposition}\label{prop:app-br-congruence}
If $\mathcal X\xrightarrow{x}\mathcal Y$ and $\mathcal X\xrightarrow{y}\mathcal Y$, then
$x\equiv y$. Consequently $L_n$ is Clark--Eyraud substitutable.
\end{proposition}

\begin{proof}
Induct on $|x|+|y|$. If $x$ or $y$ is empty, then $\mathcal X=\mathcal Y$, and the
other is empty by (P2). If $x=cx'$ and $y=cy'$ for a letter $c$,
determinism gives $\mathcal X\xrightarrow{c}\mathcal X'$ with
$\mathcal X'\xrightarrow{x'}\mathcal Y$ and $\mathcal X'\xrightarrow{y'}\mathcal Y$, and
induction applies. Otherwise Lemma~\ref{lem:app-br-branch} gives
$x=wx'$, $y=w'y'$; induction gives $x'\equiv y'$,
Lemma~\ref{lem:app-br-exchange} gives $w\equiv w'$, and $\equiv$ is a
congruence.

For the consequence, let $x,y\in\Gamma_{\mathrm{br}}^+$ and
$uxv,uyv\in L_n$. Lemma~\ref{lem:app-br-recognition} gives runs
$E_n\xrightarrow{u}\mathcal X\xrightarrow{x}\mathcal Y\xrightarrow{v}\eps$ and
$E_n\xrightarrow{u}\mathcal X'\xrightarrow{y}\mathcal Y'\xrightarrow{v}\eps$. By
(P1), $\mathcal X=\mathcal X'$ and $\mathcal Y=\mathcal Y'$, so $x\equiv y$.
\end{proof}

%==============================================================================
% Introduction: L_R example (its proof was in the comparison appendix)
%==============================================================================
For instance, one fixed window--group typing places
every language
\[
 L_R=\{\mathtt a^m\mathtt p_i\mathtt b^n
       \mathtt c^m\mathtt q_j\mathtt d^n:
       m,n\ge1,\ (i,j)\in R\},\qquad R\subseteq\{0,1\}^2,
\]
in $\Cmcf{2}{s}$: the group component records the feature values $i,j$,
while the licensing relation $R$ and the crossing dependencies must be
learned from data.

%==============================================================================
% Introduction: remaining-question paragraph (version with Theorem C(i))
%==============================================================================
\paragraph{The remaining question.}
Theorem~\ref{thm:B} concerns pseudovarieties; the cost of individual typings
outside $\mathbf{LG}$ depends on the typing.  The binary typing $s_{U,2}$ admits
polynomial data at fan-out two (Theorem~\ref{thm:C}(i)); whether it does at
fan-out one, and whether typings onto $U_1$ with two letters of type $1$ do at
fan-out two, is open (Section~\ref{sec:conclusion}).  A typing that needs
superpolynomial data at fan-out one but polynomial data at fan-out two would
show that raising the fan-out can make learning strictly cheaper.
Theorem~\ref{thm:C} shows that the fan-out-one obstruction does not lift, but
not yet that the boundary itself moves.

%==============================================================================
% Conclusion: typings-onto-U1 paragraph and Open Problems (version with Theorem C(i))
%==============================================================================
\paragraph{Typings onto $U_1$.}
Theorem~\ref{thm:B} is a statement about pseudovarieties.  Every finite
semigroup outside $\mathbf{LG}$ contains a copy of $U_1$
(Theorem~\ref{thm:lg-characterization}), so the cost of individual typings
outside $\mathbf{LG}$ is decided first for typings onto $U_1$, and there the
number of letters of type $1$ appears to matter.  Partial results, to be
reported separately, are as follows.  At fan-out one, a typing onto $U_1$ with
two letters of type $1$ and one letter of type $\zeta$ needs exponential data,
by a three-letter variant of the bracketed doubling family.  The binary typing
$s_{U,2}$ needs data quadratic in $\bar\theta_G$ for presentations of size
$O(\log\bar\theta_G)$, and this is sharp for targets on which
$|w|_{\mathtt a}+\lambda|w|_{\mathtt b}$ is constant.  Two questions remain.

\begin{openproblem}[The binary typing at fan-out one]
\label{op:S1}
For $s_{U,2}$, fan-out one and $r_{\max}\ge2$: is the characteristic-data length
bounded by a polynomial in $\|G\|$ and $\bar\theta_G$?  Two words
$\mathtt a^p,\mathtt a^q$ ($p<q$) of a target share the context $(\eps,\eps)$ and
force $\mathtt a^{p+k(q-p)}\in L$ for all $k\ge0$, so a finite target has at most
one word of type $1$.  A negative answer, together with
Theorem~\ref{thm:binary-u1-polydata}, would show that one typing needs
superpolynomial data at fan-out one and polynomial data at fan-out two.
\end{openproblem}

\begin{openproblem}[Typings outside local groups at higher fan-out]
\label{op:fanout-two}
Fix $f\ge2$ and $r_{\max}\ge2$.  For which typings $s$ with
$\mathsf S\notin\mathbf{LG}$ is there a polynomial $p$ such that, for every
$L\in\Cmcf{f}{s}$ and every reduced good fan-out-$\le f$, rank-$\le r_{\max}$
presentation $G$ of $L\cap\Sigma^+$, the operator
$\mathcal B^{\le r_{\max}}_{f,s}$ has a characteristic sample of total length
at most $p(\|G\|,\bar\theta_G)$?  At $f=2$ the typing $s_{U,2}$ is such a typing
(Theorem~\ref{thm:binary-u1-polydata}); typings onto $U_1$ with two letters of
type $1$, and all $f\ge3$, are open.  A superpolynomial lower bound for one
typing onto $U_1$ at fan-out $f$ would make $\mathbf{LG}$ the boundary at that
fan-out as well; by Proposition~\ref{prop:rare-type}, a family with
$\mu^\star(L_n,s)=n^{\omega(1)}$ and $\|G_n\|\bar\theta_{G_n}=n^{O(1)}$ would
suffice, and Theorem~\ref{thm:doubling-no-lift} excludes the bracketed doubling
family.  Conversely, polynomial data at $f=2$ for a typing onto $U_1$ with two
letters of type $1$ would show that one typing needs exponential data at fan-out
one and polynomial data at fan-out two.
\end{openproblem}

%==============================================================================
% Appendix C: full proofs of Theorem poly and Corollary gold-explanatory
%==============================================================================
\begin{proof}[Proof of Theorem~\ref{thm:poly}]
Put $n_K=\max\{1,\|K\|_+\}$.  An occurrence is a sample word with at most $2f$ cut
positions, so there are $n_K^{O(f)}$ of them, each typed in linear time; grouping
occurrences by context and type lists the Link rules.  A composition witness of rank
$\le r_{\max}$ is a parent occurrence with at most $2fr_{\max}$ further cut positions
and a labelling whose number depends only on $f,r_{\max}$; goodness is checked
directly.  This gives $n_K^{O(fr_{\max})}$ Comp rules, and all other rules are fewer.
\end{proof}

\begin{proof}[Proof of Corollary~\ref{cor:gold-explanatory}]
Let $L\in\Cmcf{f}{s}$ (for (ii), also $L\in\mathsf{MCFL}(f,r_{\max})$), let $\mathcal F_L$
be the characteristic sample of Theorem~\ref{thm:exact-reconstruction} for the
operator used, and consider a text for $L$; membership in MCFLs is decidable
\cite{SekiEtAl1991}, so the learner is computable.  Let $n_0$ be a stage by which all
of $\mathcal F_L$ has appeared.  If the learner changes its hypothesis at some stage
$n\ge n_0$, the new hypothesis $\mathcal B(K_n)$ generates exactly $L$, so it is never
changed again.  Otherwise the hypothesis $H$ held at $n_0$ is final.  If $H=H_0$, every
datum was generated by $\emptyset$, so $L=\emptyset$.  Else $H=\mathcal B(K_m)$ for some
$m\le n_0$, so $L(H)\subseteq L$ by Proposition~\ref{prop:reconstruction-soundness};
every word of $L$ appears in the text, before $m$ (then it lies in $K_m\subseteq L(H)$
by sample consistency) or after it (then it is generated by $H$, which is never
changed).  So $L(H)=L$.  For (ii), each update is a construction
(Theorem~\ref{thm:poly}) plus a membership test, both polynomial in $\|K\|_+$ for fixed
$f$ and $r_{\max}$ \cite[Proposition~4]{Yoshinaka2011}.
\end{proof}

These bounds concern the work per datum.  They do not bound the stage of convergence,
since a text may postpone any example arbitrarily long; short characteristic data alone
gives no presentation-independent convergence stage.
